% EGE 2027, task 14. Source IDs and URLs are kept in comments.
% Editable task bank. Append new BankTask blocks at the end of the document.
\documentclass[12pt,a4paper]{article}
\usepackage{physics-textbook}
\usepackage{graphicx,needspace}
\newif\ifShowSolutions
\ifdefined\TutorEdition
  \ShowSolutionstrue
\else
  \ShowSolutionsfalse
\fi
\newcounter{banktask}
\newcommand{\BankTask}[3]{%
  \refstepcounter{banktask}%
  \Needspace{35mm}\par\medskip\noindent\begin{minipage}{\linewidth}%
  \noindent\textbf{\thebanktask.}\quad #2\par
  \ifShowSolutions
    #3\par
    \ifdefined\TutorEdition
      \ifcsname TutorSource#1\endcsname
        \smallskip{\small\csname TutorSource#1\endcsname\par}%
      \fi
    \fi
  \fi
  \end{minipage}\par\medskip
}
\begin{document}


% Source 7D9e2D: https://ege.fipi.ru/bank/questions.php?search=1&proj=BA1F39653304A5B041B656915DC36B38&qid=7D9e2D
\BankTask{7D9e2D}{
Две маленькие закреплённые проводящие бусинки, расположенные в точках $A$ и $B$, несут на себе заряды $+3q>0$ и $-1{,}5q<0$ соответственно (см. рисунок). Точка $C$ находится посередине между бусинками $A$ и $B$.
\begin{enumerate}
\item На бусинку $B$ со стороны бусинки $A$ действует сила Кулона, направленная горизонтально вправо.
\item Напряжённость результирующего электростатического поля в точке $C$ направлена горизонтально влево.
\item Модули сил Кулона, действующих на бусинки, одинаковы.
\item Если бусинки соединить тонкой стальной проволокой, они начнут отталкиваться друг от друга.
\item Если бусинки соединить незаряженной стеклянной палочкой, их заряды станут одинаковыми.
\end{enumerate}
\begin{center}
\begin{tikzpicture}[book diagram]
\draw[densely dashed] (0,0)--(4,0);
\fill (0,0) circle (2pt) node[above=7pt] {$+3q$} node[below=5pt] {$A$};
\fill (2,0) circle (2pt) node[below=5pt] {$C$};
\fill (4,0) circle (2pt) node[above=7pt] {$-1{,}5q$} node[below=5pt] {$B$};
\end{tikzpicture}
\end{center}
}{
\smallskip\noindent\textbf{Решение от GPT}\par
1) Неверно: разноимённые заряды притягиваются, поэтому сила на бусинку $B$ направлена влево. 2) Неверно: в точке $C$ оба поля направлены вправо — от положительного заряда к отрицательному. 3) Верно по третьему закону Ньютона: $F_{AB}=F_{BA}$. 4) Верно: при соединении проводником суммарный положительный заряд $3q-1{,}5q=1{,}5q$ перераспределится, и бусинки будут отталкиваться. 5) Неверно: стекло — диэлектрик и не обеспечивает перетекания зарядов между бусинками. Ответ: $3,\ 4$.
}

% Source E5F230: https://ege.fipi.ru/bank/questions.php?search=1&proj=BA1F39653304A5B041B656915DC36B38&qid=E5F230
\BankTask{E5F230}{
Однородное электростатическое поле создано равномерно заряженной протяженной горизонтальной пластиной. Линии напряжённости поля направлены вертикально вверх (см. рисунок).
\begin{enumerate}
\item Пластина имеет отрицательный заряд.
\item Потенциал электростатического поля в точке $A$ ниже, чем в точке $C$.
\item Работа электростатического поля по перемещению пробного точечного отрицательного заряда из точки $A$ в точку $B$ отрицательна.
\item Если в точку $B$ поместить пробный точечный положительный заряд, то на него со стороны пластины будет действовать сила, направленная вертикально вверх.
\item Напряжённость поля в точках $A$, $B$ и $C$ одинакова.
\end{enumerate}
\begin{center}
\begin{tikzpicture}[book diagram]
\draw[line width=1.2pt] (-0.45,0)--(2.05,0);
\foreach \x in {0,0.55,1.1,1.65}{
\draw (\x,0)--(\x,2.55);
\draw[-{book tip}] (\x,1.65)--(\x,2.03);
}
\draw[dashed] (0,0.8)--(1.65,0.8);
\draw[dashed] (0,0.8)--(1.65,1.68);
\draw (1.50,0.8)--(1.50,0.95)--(1.65,0.95);
\fill (0,0.8) circle (2pt) node[left=3pt] {$A$};
\fill (1.65,0.8) circle (2pt) node[right=3pt] {$B$};
\fill (1.65,1.68) circle (2pt) node[right=3pt] {$C$};
\node at (1.35,2.48) {$\vec E$};
\end{tikzpicture}
\end{center}
}{
\smallskip\noindent\textbf{Решение от GPT}\par
1) Неверно: линии поля направлены от пластины, следовательно, её заряд положителен. 2) Неверно: потенциал убывает вдоль поля, поэтому $\varphi_A>\varphi_C$. 3) Неверно: точки $A$ и $B$ лежат на одной эквипотенциальной поверхности, так что $A_{AB}=q(\varphi_A-\varphi_B)=0$. 4) Верно: для положительного заряда сила $\vec F=q\vec E$ направлена вверх. 5) Верно: в однородном поле $\vec E_A=\vec E_B=\vec E_C$. Ответ: $4,\ 5$.
}

% Source B6C47E: https://ege.fipi.ru/bank/questions.php?search=1&proj=BA1F39653304A5B041B656915DC36B38&qid=B6C47E
\BankTask{B6C47E}{
На уединённой неподвижной проводящей сфере радиусом $R$ находится положительный заряд $Q$. Сфера находится в вакууме. Напряжённость электростатического поля сферы в точке $A$ равна $36\ \text{В/м}$. Все расстояния указаны на рисунке.
\begin{enumerate}
\item Напряжённость поля в точке $B$ $E_B=576\ \text{В/м}$.
\item Напряжённость поля в точке $C$ $E_C=36\ \text{В/м}$.
\item Потенциал электростатического поля в точке $B$ выше, чем в точке $D$: $\varphi_B>\varphi_D$.
\item Потенциал электростатического поля в точках $D$ и $F$ одинаков: $\varphi_D=\varphi_F$.
\item Потенциал электростатического поля в точке $C$ выше, чем в точке $F$: $\varphi_C>\varphi_F$.
\end{enumerate}
\begin{center}
\begin{tikzpicture}[book diagram]
\draw (0,0) circle (1.45);
\draw[{book tip}-{book tip}] (-0.725,0)--(0,0) node[midway,above=5pt] {$\dfrac{R}{2}$};
\draw[{book tip}-{book tip}] (0,0)--(1.45,0) node[midway,above=3pt] {$R$};
\draw[{book tip}-{book tip}] (1.45,0)--(2.9,0) node[midway,above=3pt] {$R$};
\draw[{book tip}-{book tip}] (0,0)--(0,-1.45) node[midway,left=3pt] {$R$};
\draw[{book tip}-{book tip}] (0,-1.45)--(0,-2.9) node[midway,left=3pt] {$R$};
\fill (-0.725,0) circle (1.6pt) node[above left=2pt] {$B$};
\fill (0,0) circle (1.6pt);
\fill (1.45,0) circle (1.6pt) node[below right=2pt] {$D$};
\fill (2.9,0) circle (1.6pt) node[above right=1pt] {$A$};
\fill (0,-1.45) circle (1.6pt) node[below right=2pt] {$F$};
\fill (0,-2.9) circle (1.6pt) node[below=3pt] {$C$};
\end{tikzpicture}
\end{center}
}{
\smallskip\noindent\textbf{Решение от GPT}\par
1) Неверно: внутри проводящей сферы электростатическое поле отсутствует, поэтому $E_B=0$. 2) Верно: $r_A=r_C=2R$, поэтому $E_C=E_A=kQ/(2R)^2=36\ \text{В/м}$. 3) Неверно: внутри проводника и на его поверхности потенциал одинаков, $\varphi_B=\varphi_D=kQ/R$. 4) Верно: точки $D$ и $F$ находятся на поверхности сферы, поэтому $\varphi_D=\varphi_F$. 5) Неверно: $\varphi_C=kQ/(2R)<kQ/R=\varphi_F$. Ответ: $2,\ 4$.
}

% Source 896222: https://ege.fipi.ru/bank/questions.php?search=1&proj=BA1F39653304A5B041B656915DC36B38&qid=896222
\BankTask{896222}{
Незаряженное металлическое тело, продольное сечение которого показано на рисунке, поместили в однородное электрическое поле напряжённостью $\vec{E}$.

Какие два утверждения описывают результаты воздействия этого поля на металлическое тело?
\begin{enumerate}
\item Напряжённость электрического поля в точке $D$ равна нулю.
\item Потенциалы в точках $B$ и $C$ равны.
\item Концентрация свободных электронов в точке $A$ наибольшая.
\item В точке $A$ индуцируется положительный заряд.
\item В точке $D$ индуцируется отрицательный заряд.
\end{enumerate}
\begin{center}
\begin{tikzpicture}[book diagram,x=1cm,y=1cm]

\draw[-{book tip}] (0,2.45)--(3.8,2.45);
\draw[-{book tip}] (0,0.15)--(3.8,0.15);
\node[above] at (3.6,2.55) {$\vec E$};
\node[above] at (3.6,0.25) {$\vec E$};
\draw[fill=gray!55,line width=0.6pt] (0.85,0.87)--(3.05,0.87) arc[start angle=-90,end angle=90,radius=0.45]--(0.85,1.77) arc[start angle=90,end angle=270,radius=0.45]--cycle;
\fill (1.95,1.77) circle (0.035);
\fill (0.4,1.32) circle (0.035);
\fill (3.5,1.32) circle (0.035);
\fill (1.95,1.36) circle (0.04);
\node[above] at (1.95,1.82) {$A$};
\node[left] at (0.4,1.32) {$B$};
\node[right] at (3.5,1.32) {$C$};
\node[right] at (1.97,1.3) {$D$};
\end{tikzpicture}
\end{center}
}{
\smallskip\noindent\textbf{Решение от GPT}\par
В электростатическом равновесии внутри металла $\vec E=0$, а весь проводник эквипотенциален: $\varphi_B=\varphi_C$; утверждения 1 и 2 верны. Электроны смещаются против внешнего поля и накапливаются у левого края $B$, поэтому утверждение 3 неверно. В точке $A$ на середине боковой поверхности положительный заряд не индуцируется, а внутри металла в точке $D$ избыточного заряда нет, поэтому утверждения 4 и 5 неверны. Ответ: $12$.
}

% Source E412A6: https://ege.fipi.ru/bank/questions.php?search=1&proj=BA1F39653304A5B041B656915DC36B38&qid=E412A6
\BankTask{E412A6}{
Стеклянные незаряженные кубики 1 и 2 и медные незаряженные кубики 3 и 4 сблизили вплотную и поместили в электрическое поле, напряжённость которого направлена горизонтально вправо, как показано в левой части рисунка. Затем кубики раздвинули и уже после этого выключили электрическое поле (правая часть рисунка).

Какие два утверждения описывают данный процесс?
\begin{enumerate}
\item В электрическом поле кубики 1 и 2 приобретают суммарный положительный заряд.
\item При помещении стеклянных кубиков в электрическое поле наблюдается явление поляризации.
\item В электрическом поле кубики 3 и 4 приобретают суммарный отрицательный заряд.
\item После разделения кубик 1 имеет положительный заряд.
\item После разделения кубик 4 имеет положительный заряд.
\end{enumerate}
\begin{center}
\begin{tikzpicture}[book diagram,x=1cm,y=1cm]
\foreach \x/\y in {1.3/1.25,2.55/1.25,5.25/1.25,7.1/1.25}{
\fill[gray!8] (\x,\y) rectangle ++(1.25,1.25);
\foreach \h in {0.06,0.12,...,1.20}{\draw[gray!55,line width=0.2pt] (\x,{\y+\h})--++(1.25,0);}
\draw (\x,\y) rectangle ++(1.25,1.25);
}
\foreach \x/\y in {1.3/0,2.55/0,5.3/-0.75,7.1/-0.75}{
\fill[gray!30] (\x,\y) rectangle ++(1.25,1.25);
\foreach \h in {0.06,0.12,...,1.20}{\draw[gray!45,line width=0.2pt] (\x,{\y+\h})--++(1.25,0);}
\draw (\x,\y) rectangle ++(1.25,1.25);
}
\node at (1.925,1.875) {1};
\node at (3.175,1.875) {2};
\node at (1.925,0.625) {3};
\node at (3.175,0.625) {4};
\node at (5.875,1.875) {1};
\node at (7.725,1.875) {2};
\node at (5.925,-0.125) {3};
\node at (7.725,-0.125) {4};
\draw[-{book tip}] (-0.1,1.875)--(1.2,1.875);
\draw[-{book tip}] (-0.1,0.625)--(1.2,0.625);
\node[above] at (0.55,1.95) {$\vec E$};
\node[above] at (0.55,0.7) {$\vec E$};
\draw[dashed,gray] (4.48,-1.1)--(4.48,2.9);
\draw[dashed,gray] (4.72,-1.1)--(4.72,2.9);
\end{tikzpicture}
\end{center}
}{
\smallskip\noindent\textbf{Решение от GPT}\par
Стекло поляризуется, но суммарный заряд каждого стеклянного кубика остаётся нулевым: утверждение 2 верно, а 1 и 4 неверны. Общий заряд медных кубиков сохраняется: $q_3+q_4=0$, поэтому утверждение 3 неверно. Электроны переходят из кубика 4 в кубик 3 против направления поля; после разделения $q_3<0$, $q_4>0$, поэтому утверждение 5 верно. Ответ: 2, 5.
}

% Source 2786AC: https://ege.fipi.ru/bank/questions.php?search=1&proj=BA1F39653304A5B041B656915DC36B38&qid=2786AC
\BankTask{2786AC}{
В устройстве, изображённом на рисунке, лёгкий шарик из алюминиевой фольги подвешен между пластинами плоского конденсатора на изолирующей нити, при этом одна из пластин заземлена, а другая заряжена положительно. Когда устройство собрано, шарик приходит в колебательное движение, при котором сталкивается с пластинами конденсатора.

Какие утверждения соответствуют колебательному движению шарика в конденсаторе после первого столкновения шарика с пластиной?
\begin{enumerate}
\item Пренебрегая сопротивлением воздуха при движении шарика, его колебания можно считать гармоническими.
\item При движении шарика к заземлённой пластине он заряжен положительно, а при движении к положительно заряженной пластине -- отрицательно.
\item При движении шарика к положительно заряженной пластине его заряд равен нулю, а при движении к заземлённой пластине его заряд -- положительный.
\item По мере колебаний шарика электрическая ёмкость конденсатора увеличивается.
\item По мере колебаний шарика заряд конденсатора уменьшается.
\end{enumerate}
\begin{center}
\begin{tikzpicture}[book diagram,x=1cm,y=1cm,line cap=round,line join=round]
% Верхняя опора и подвес.
\fill[gray!60] (0.9,6.2) rectangle (3.45,6.32);
\draw[line width=1pt] (0.9,6.32)--(3.45,6.32);
\draw (2.18,6.2)--(2.18,3.16);
\shade[ball color=gray!25] (2.18,3.04) circle (0.16);
\draw (2.18,3.04) circle (0.16);
% Основание установки.
\draw[fill=gray!45] (0,0.72)--(3.95,0.72)--(4.06,0.9)--(3.95,1.08)--(0,1.08)--(0.1,0.9)--cycle;
% Изогнутые стойки.
\draw[fill=gray!55] (0.73,1.08)--(0.73,2.91) .. controls (0.73,3.12) and (0.87,3.19) .. (1.08,3.19)--(1.35,3.19)--(1.35,2.98)--(1.1,2.98) .. controls (1,2.98) and (0.95,2.94) .. (0.95,2.84)--(0.95,1.08)--cycle;
\draw[fill=gray!55] (3.48,1.08)--(3.48,2.91) .. controls (3.48,3.12) and (3.34,3.19) .. (3.13,3.19)--(2.88,3.19)--(2.88,2.98)--(3.11,2.98) .. controls (3.21,2.98) and (3.26,2.94) .. (3.26,2.84)--(3.26,1.08)--cycle;
% Крепления пластин.
\draw[fill=gray!75] (1.25,2.91) rectangle (1.49,3.25);
\draw[fill=gray!75] (2.88,2.91) rectangle (3.13,3.25);
% Пластины конденсатора.
\draw[fill=gray!15] (1.49,2.01) rectangle (1.55,4.04);
\draw[fill=gray!15] (2.83,2.01) rectangle (2.89,4.04);
\node at (1.28,3.94) {$+$};
% Зажимы на основании.
\foreach \x in {0.84,3.37}{
\draw[fill=gray!70] ({\x-0.23},0.64) rectangle ({\x+0.23},1.13);
\draw[fill=gray!65] ({\x-0.15},1.13) rectangle ({\x+0.15},1.32);
\draw[fill=black] ({\x-0.065},0.64)--({\x+0.065},0.64)--({\x+0.065},0.43)--({\x+0.17},0.35) .. controls ({\x+0.2},0.25) and ({\x-0.2},0.25) .. ({\x-0.17},0.35)--({\x-0.065},0.43)--cycle;
}
% Провод к земле.
\draw (2.89,3.8) .. controls (3.03,3.82) and (3.02,3.7) .. (3.2,3.72) .. controls (3.42,3.67) and (3.55,3.65) .. (3.65,3.47) .. controls (3.83,3.44) and (3.79,3.27) .. (3.93,3.18) .. controls (4.12,2.97) and (4.14,2.8) .. (4.28,2.64) .. controls (4.45,2.39) and (4.3,2.25) .. (4.4,2.12) .. controls (4.54,1.94) and (4.43,1.83) .. (4.53,1.68) .. controls (4.63,1.53) and (4.5,1.35) .. (4.61,1.2) .. controls (4.78,0.96) and (4.73,0.82) .. (4.87,0.59) .. controls (5.02,0.36) and (4.91,0.19) .. (4.91,0.03);
\draw (4.68,0.03)--(5.15,0.03);
\foreach \x in {4.7,4.78,4.86,4.94,5.02,5.1}{\draw (\x,0.03)--++(-0.1,-0.15);}
\end{tikzpicture}
\end{center}
}{
\smallskip\noindent\textbf{Решение от GPT}\par
При контакте с положительной пластиной шарик получает положительный заряд, а при контакте с заземлённой в поле конденсатора — отрицательный: утверждение 2 верно, а 3 неверно. Сила $\vec F=q\vec E$ не является возвращающей силой вида $F_x=-kx$, а столкновения меняют заряд шарика, поэтому колебания не гармонические и утверждение 1 неверно. Ёмкость $C=\varepsilon_0 S/d$ определяется геометрией конденсатора и не увеличивается: утверждение 4 неверно. Шарик переносит заряд между пластинами, разряжая конденсатор, поэтому утверждение 5 верно. Ответ: 2, 5.
}

% Source 0cDA83: https://ege.fipi.ru/bank/questions.php?search=1&proj=BA1F39653304A5B041B656915DC36B38&qid=0cDA83
\BankTask{0cDA83}{
Плоский воздушный конденсатор ёмкостью \(C_0\) состоит из двух металлических пластин, находящихся на расстоянии \(d_0\) друг от друга. Конденсатор зарядили и отключили от источника постоянного напряжения. Расстояние между пластинами меняется со временем так, как показано на графике.
\begin{enumerate}
\item В интервале времени от \(0\) до \(t_4\) ёмкость конденсатора остаётся неизменной.
\item В интервале времени от \(t_1\) до \(t_4\) заряд конденсатора остаётся неизменным.
\item В интервале времени от \(0\) до \(t_4\) энергия конденсатора уменьшается.
\item В интервале времени от \(0\) до \(t_4\) напряжённость электрического поля между пластинами конденсатора остаётся постоянной.
\item В интервале времени от \(0\) до \(t_4\) напряжение между пластинами конденсатора уменьшается.
\end{enumerate}
\begin{center}
\begin{tikzpicture}[book diagram,x=0.65cm,y=0.65cm]
\foreach \y in {1,...,5}{\draw[thin] (0.35,\y)--(2.35,\y);}
\foreach \x in {1,2}{\draw[thin] (\x,1.33)--(\x,4.67);}
\foreach \x in {3,4}{\draw[dashed,thin] (\x,0)--(\x,5);}
\draw[-{book tip}] (0,0)--(5.2,0) node[below] {$t$};
\draw[-{book tip}] (0,0)--(0,6) node[left] {$d$};
\node[below left] at (0,0) {$0$};
\foreach \i in {1,...,4}{\draw (\i,-0.12)--(\i,0.12);\node[below] at (\i,-0.12) {$t_{\i}$};}
\foreach \j in {1,...,5}{\draw (-0.12,\j)--(0.12,\j);}
\node[left] at (-0.12,1) {$d_0$};
\draw[very thick] (0,1)--(4,4);
\end{tikzpicture}
\end{center}
}{
\smallskip\noindent\textbf{Решение от GPT}\par
После отключения от источника заряд \(Q\) сохраняется, поэтому утверждение 2 верно. При увеличении \(d\) ёмкость \(C=\varepsilon_0S/d\) уменьшается, а энергия \(W=Q^2/(2C)\) увеличивается: утверждения 1 и 3 неверны. Поле \(E=Q/(\varepsilon_0S)\) постоянно, поэтому утверждение 4 верно. Напряжение \(U=Ed\) увеличивается, поэтому утверждение 5 неверно. Ответ: \(2,\ 4\).
}

% Source 675443: https://ege.fipi.ru/bank/questions.php?search=1&proj=BA1F39653304A5B041B656915DC36B38&qid=675443
\BankTask{675443}{
Плоский воздушный конденсатор ёмкостью \(C_0\), подключённый к источнику постоянного напряжения, состоит из двух металлических пластин, находящихся на расстоянии \(d_0\) друг от друга. Расстояние между пластинами меняется со временем так, как показано на графике.
\begin{enumerate}
\item В интервале времени от \(t_1\) до \(t_4\) заряд конденсатора возрастает.
\item В интервале времени от \(t_1\) до \(t_4\) энергия конденсатора равномерно уменьшается.
\item В промежутке времени от \(t_1\) до \(t_4\) напряжённость электрического поля между пластинами конденсатора остаётся постоянной.
\item В момент времени \(t_4\) ёмкость конденсатора увеличилась в $5$ раз по сравнению с первоначальной (при \(t=0\)).
\item В промежутке времени от \(t_1\) до \(t_4\) напряжённость электрического поля между пластинами конденсатора убывает.
\end{enumerate}
\begin{center}
\begin{tikzpicture}[book diagram,x=0.9cm,y=0.6cm]
\draw[thin] (0,0) grid[step=1] (4,6);
\draw[-{book tip}] (0,0)--(4.7,0) node[right] {$t$};
\draw[-{book tip}] (0,0)--(0,6.7) node[above] {$d$};
\node[below left] at (0,0) {$0$};
\node[left] at (0,5) {$d_0$};
\foreach \i in {1,...,4}{\node[below] at (\i,0) {$t_{\i}$};}
\draw[thick] (0,5)--(4,1);
\end{tikzpicture}
\end{center}
}{
\smallskip\noindent\textbf{Решение от GPT}\par
При постоянном напряжении $U$ расстояние $d$ уменьшается, поэтому $C=\varepsilon_0S/d$, $Q=CU$, $W=CU^2/2$ и $E=U/d$ возрастают. Утверждение 1 верно, утверждения 2, 3 и 5 неверны. По графику $d(t_4)=d_0/5$, следовательно, $C(t_4)=5C_0$: утверждение 4 верно. Ответ: $14$.
}

% Source 410957: https://ege.fipi.ru/bank/questions.php?search=1&proj=BA1F39653304A5B041B656915DC36B38&qid=410957
\BankTask{410957}{
Две параллельные металлические пластины больших размеров расположены на расстоянии \(d\) друг от друга и подключены к источнику постоянного напряжения (рис. 1). Пластины закрепили на изолирующих подставках и спустя длительное время отключили от источника (рис. 2).
\begin{enumerate}
\item Напряжённость электрического поля в точке \(B\) больше, чем в точке \(C\).
\item Потенциалы электрического поля в точках \(A\) и \(B\) одинаковы.
\item Если после отключения от источника увеличить расстояние \(d\) между пластинами, то напряжённость электрического поля в точке \(A\) увеличится.
\item Если после отключения от источника уменьшить расстояние \(d\) между пластинами, то заряд правой пластины не изменится.
\item Если после отключения от источника пластины полностью погрузить в керосин, то энергия электрического поля системы пластин уменьшится.
\end{enumerate}
\begin{center}
\begin{tikzpicture}[book diagram,x=0.8cm,y=0.8cm]
\draw[very thick] (0,0.7)--(0,3.5) (1.25,0.7)--(1.25,3.5);
\draw[thin] (0,3.5)--(0,4.35) (1.25,3.5)--(1.25,4.35);
\draw[{book tip}-{book tip}] (0,3.9)--(1.25,3.9) node[midway,above] {$d$};
\draw (0,1.85)--(-1.1,1.85)--(-1.1,-0.2)--(0.35,-0.2);
\draw (1.25,1.85)--(2.35,1.85)--(2.35,-0.2)--(0.9,-0.2);
\draw[fill=white] (0.35,-0.2) circle (0.065) (0.9,-0.2) circle (0.065);
\node[above] at (0.35,-0.15) {$-$};
\node[above] at (0.9,-0.15) {$+$};
\node at (0.625,-0.8) {Рис. 1};
\begin{scope}[xshift=4.7cm]
\draw[very thick] (0,0.7)--(0,3.5) (1.25,0.7)--(1.25,3.5);
\draw[thin] (0,3.5)--(0,4.35) (1.25,3.5)--(1.25,4.35);
\draw[{book tip}-{book tip}] (0,3.9)--(1.25,3.9) node[midway,above] {$d$};
\draw (0,1.85)--(-0.9,1.85) (1.25,1.85)--(2.15,1.85);
\fill (0.35,2.35) circle (0.045) node[above left] {$B$};
\fill (0.98,1.85) circle (0.045) node[above left] {$C$};
\fill (0.65,1.15) circle (0.045) node[left] {$A$};

\node at (0.625,-0.8) {Рис. 2};
\end{scope}
\end{tikzpicture}
\end{center}
}{
\smallskip\noindent\textbf{Решение от GPT}\par
Между большими параллельными пластинами поле однородно, поэтому \(E_B=E_C\) и утверждение 1 неверно. Точки \(A\) и \(B\) находятся на разных расстояниях от пластин, поэтому их потенциалы различаются: утверждение 2 неверно. После отключения заряд сохраняется, а \(E=Q/(\varepsilon_0S)\) не зависит от \(d\), поэтому утверждение 3 неверно, а 4 верно. При погружении в керосин ёмкость увеличивается в \(\varepsilon\) раз, а энергия \(W=Q^2/(2C)\) уменьшается в \(\varepsilon\) раз, поэтому утверждение 5 верно. Ответ: \(4,\ 5\).
}

% Source 5D79D1: https://ege.fipi.ru/bank/questions.php?search=1&proj=BA1F39653304A5B041B656915DC36B38&qid=5D79D1
\BankTask{5D79D1}{
При нагревании спирали лампы накаливания протекающим по ней электрическим током основная часть подводимой энергии теряется в виде теплового излучения. На рисунке изображены графики зависимости мощности тепловых потерь лампы от температуры спирали $P=P(T)$ и силы тока от приложенного напряжения $I=I(U)$.
\begin{enumerate}
\item С уменьшением приложенного напряжения температура спирали уменьшается пропорционально мощности.
\item При напряжении $100\,\text{В}$ сопротивление спирали лампы равно $50\,\text{Ом}$.
\item Сопротивление спирали лампы постоянно и не зависит от температуры.
\item Температура спирали лампы при приложенном напряжении $U=120\,\text{В}$ меньше $3600\,\text{К}$.
\item С увеличением силы протекающего тока температура спирали увеличивается.
\end{enumerate}
\begin{center}
\begin{tikzpicture}[book diagram,x=1cm,y=1cm]
\begin{scope}
\draw[thin,step=0.5] (0,0) grid (4,2.5);
\draw[-{book tip}] (0,0)--(4.85,0) node[below] {$T,\,\text{К}$};
\draw[-{book tip}] (0,0)--(0,2.95) node[left] {$P,\,\text{Вт}$};
\foreach \x/\t in {0/2000,1/2400,2/2800,3/3200,4/3600}{\node[below,font=\small] at (\x,0) {$\t$};}
\foreach \y/\p in {0/0,1/100,2/200}{\node[left,font=\small] at (0,\y) {$\p$};}
\draw[thin] plot[smooth] coordinates {(0,0.17) (0.5,0.25) (1,0.40) (1.5,0.58) (2,0.77) (2.5,1.02) (3,1.33) (3.5,1.71) (4,2.14)};
\end{scope}
\begin{scope}[xshift=6.3cm]
\draw[thin,step=0.4] (0,0) grid (4,2.8);
\draw[-{book tip}] (0,0)--(4.85,0) node[below] {$U,\,\text{В}$};
\draw[-{book tip}] (0,0)--(0,3.15) node[left] {$I,\,\text{А}$};
\foreach \x/\u in {0/50,2/100,4/150}{\node[below,font=\small] at (\x,0) {$\u$};}
\foreach \y/\i in {0/{1{,}2},0.8/{1{,}6},1.6/{2{,}0},2.4/{2{,}4}}{\node[left,font=\small] at (0,\y) {$\i$};}
\draw[thin] plot[smooth] coordinates {(0,0.4) (0.4,0.65) (0.8,0.89) (1.2,1.13) (1.6,1.37) (2,1.6) (2.4,1.82) (2.8,2.03) (3.2,2.22) (3.6,2.40) (4,2.56)};
\end{scope}
\end{tikzpicture}
\end{center}
}{
\smallskip\noindent\textbf{Решение от GPT}\par
Утверждение 1 неверно: график $P(T)$ нелинейный, поэтому температура не пропорциональна мощности. Утверждение 2 верно: при $U=100\,\text{В}$ ток $I=2\,\text{А}$, следовательно, $R=U/I=50\,\text{Ом}$. Утверждение 3 неверно: отношение $U/I$ меняется, например при $50\,\text{В}$ оно равно примерно $36\,\text{Ом}$. Утверждение 4 неверно: при $120\,\text{В}$ мощность $UI\approx264\,\text{Вт}$ превышает тепловые потери при $3600\,\text{К}$, поэтому температура выше $3600\,\text{К}$. Утверждение 5 верно: рост тока сопровождается ростом напряжения и мощности, а значит, температуры спирали. Ответ: 2, 5.
}

% Source 56DF4A: https://ege.fipi.ru/bank/questions.php?search=1&proj=BA1F39653304A5B041B656915DC36B38&qid=56DF4A
\BankTask{56DF4A}{
Алюминиевый проводник $\text{АБ}$ подвешен на тонких медных проволочках к деревянной балке и подключён к источнику постоянного напряжения так, как показано на рисунке. Справа от проводника находится южный полюс постоянного магнита. Ползунок реостата плавно перемещают влево.
\begin{enumerate}
\item Сопротивление внешней цепи увеличивается.
\item Линии индукции магнитного поля, созданного магнитом вблизи проводника $\text{АБ}$, направлены влево.
\item Сила тока, протекающего через проводник $\text{АБ}$, уменьшается.
\item Сила Ампера, действующая на проводник $\text{АБ}$, увеличивается.
\item Сила натяжения проволочек, на которых подвешен проводник $\text{АБ}$, увеличивается.
\end{enumerate}
\begin{center}
\begin{tikzpicture}[book diagram,x=0.027cm,y=0.027cm,line cap=round,line join=round]
\draw[line width=1pt] (1,141)--(134,207);
\foreach \x in {4,14,...,114}{\draw[thin] (\x,{140+0.5*\x})--({\x-13},{153+0.5*\x});}
\draw (20,151)--(20,20);
\draw (114,187)--(114,70);
\fill (20,140) circle (1.7);
\fill (114,186) circle (1.7);
\draw[line width=2pt] (20,20)--(114,70);
\node[left] at (15,30) {$\text{А}$};
\node[right] at (114,76) {$\text{Б}$};
\draw (20,140)--(142,140)--(171,155);
\draw[line width=0.9pt] (171,133)--(171,176);
\draw[line width=0.9pt] (178,145)--(178,164);
\node at (156,132) {$+$};
\node at (193,149) {$-$};
\draw (178,155)--(236,186)--(187,186);
\draw (114,186)--(161,186);
\draw (161,180) rectangle (187,192);
\draw (198,186)--(198,207)--(171,207);
\draw[-{book tip}] (171,207)--(171,193);
\draw (107,1)--(225,1)--(265,20)--(265,58)--(143,58)--(107,39)--cycle;
\draw (107,39)--(225,39)--(265,58);
\draw (225,1)--(225,39);
\draw (167,1)--(167,39)--(202,58);
\node[font=\Large] at (122,20) {$S$};
\node[font=\Large] at (210,20) {$N$};
\end{tikzpicture}
\end{center}
}{
\smallskip\noindent\textbf{Решение от GPT}\par
При перемещении ползунка влево рабочая длина реостата и сопротивление внешней цепи уменьшаются, поэтому ток возрастает: утверждения 1 и 3 неверны. Вблизи южного полюса магнитное поле направлено к нему, то есть вправо, поэтому утверждение 2 неверно. Ток идёт по проводнику от $\text{А}$ к $\text{Б}$, и по правилу левой руки сила Ампера направлена вниз; её модуль $F_A=BIl$ возрастает, поэтому утверждение 4 верно. Для равновесия суммарная вертикальная сила натяжения равна $mg+F_A$ и также возрастает, поэтому утверждение 5 верно. Ответ: $45$.
}

% Source 4E567B: https://ege.fipi.ru/bank/questions.php?search=1&proj=BA1F39653304A5B041B656915DC36B38&qid=4E567B
\BankTask{4E567B}{
В медное кольцо № 1 вносят северный полюс полосового магнита, а в пластмассовое кольцо № 2 -- южный полюс такого же магнита (см. рисунок).
\begin{enumerate}
\item В кольце № 1 индукционный ток не возникает.
\item В кольце № 2 возникает индукционный ток.
\item Кольцо № 1 притягивается к магниту.
\item Кольцо № 2 не взаимодействует с магнитом.
\item В кольце № 1 возникает ЭДС индукции.
\end{enumerate}
\begin{center}
\begin{tikzpicture}[book diagram,x=0.025cm,y=0.025cm,line join=round]
\begin{scope}
\fill[gray!35] (48,20)--(94,20)--(113,39)--(113,80)--(67,80)--(48,61)--cycle;
\draw (1,20)--(94,20)--(113,39)--(113,80)--(21,80)--(1,61)--cycle;
\draw (1,61)--(94,61)--(113,80);
\draw (94,20)--(94,61);
\draw (48,20)--(48,61)--(67,80);
\node[font=\large] at (24,40) {$S$};
\node[font=\large] at (72,40) {$N$};
\draw[-{book tip}] (38,7)--(70,7);
\fill[gray!35] (144,53) ellipse[x radius=19,y radius=51];
\fill[white] (151,53) ellipse[x radius=12,y radius=51];
\draw (144,53) ellipse[x radius=19,y radius=51];
\draw (151,53) ellipse[x radius=12,y radius=51];
\draw (146,103) .. controls (130,84) and (131,22) .. (145,3);
\node[above] at (144,112) {№ 1};
\end{scope}
\begin{scope}[xshift=6cm]
\fill[gray!35] (1,20)--(48,20)--(48,61)--(67,80)--(21,80)--(1,61)--cycle;
\draw (1,20)--(94,20)--(113,39)--(113,80)--(21,80)--(1,61)--cycle;
\draw (1,61)--(94,61)--(113,80);
\draw (94,20)--(94,61);
\draw (48,20)--(48,61)--(67,80);
\node[font=\large] at (24,40) {$N$};
\node[font=\large] at (72,40) {$S$};
\draw[-{book tip}] (38,7)--(70,7);
\fill[gray!80] (144,53) ellipse[x radius=19,y radius=51];
\fill[white] (151,53) ellipse[x radius=12,y radius=51];
\draw (144,53) ellipse[x radius=19,y radius=51];
\draw (151,53) ellipse[x radius=12,y radius=51];
\draw (146,103) .. controls (130,84) and (131,22) .. (145,3);
\node[above] at (144,112) {№ 2};
\end{scope}
\end{tikzpicture}
\end{center}
}{
\smallskip\noindent\textbf{Решение от GPT}\par
При внесении магнита поток через медное кольцо меняется, поэтому возникает ЭДС $\mathcal E=-\Delta\Phi/\Delta t$ и индукционный ток: утверждение 1 неверно, а 5 верно. Пластмасса является диэлектриком, поэтому индукционный ток в кольце № 2 не возникает: утверждение 2 неверно. По правилу Ленца медное кольцо препятствует приближению магнита и отталкивается от него, поэтому утверждение 3 неверно. Пластмассовое кольцо в рассматриваемом опыте с магнитом не взаимодействует, поэтому утверждение 4 верно. Ответ: 4, 5.
}

% Source 063141: https://ege.fipi.ru/bank/questions.php?search=1&proj=BA1F39653304A5B041B656915DC36B38&qid=063141
\BankTask{063141}{
Катушка №~1 включена в электрическую цепь, состоящую из источника постоянного напряжения и реостата. Катушка №~2 помещена внутрь катушки №~1, и её обмотка замкнута. Вид с торца катушек представлен на рисунке.

Ползунок реостата перемещают влево.
\begin{enumerate}
\item Сила тока в катушке №~1 увеличивается.
\item Модуль вектора индукции магнитного поля, созданного катушкой №~1, увеличивается.
\item Модуль магнитного потока, пронизывающего катушку №~2, уменьшается.
\item Вектор магнитной индукции магнитного поля, созданного катушкой №~2 в её центре, направлен от наблюдателя.
\item В катушке №~2 индукционный ток направлен по часовой стрелке.
\end{enumerate}
\begin{center}
\begin{tikzpicture}[book diagram,x=0.85cm,y=0.85cm]
\draw (1.65,2.25) ++(-104:0.9) arc[start angle=-104,end angle=-436,radius=0.9];
\draw (1.65,2.25) circle[radius=0.5];
\node at (1.65,2.37) {№ 2};
\node at (0.36,1.98) {№ 1};
\draw (1.432,1.377)--(0.2,1.377)--(0.2,0.38)--(0.58,0.38);
\draw (1.868,1.377)--(3.14,1.377)--(3.14,0.38)--(2.35,0.38);
\draw (0.58,0.18) rectangle (1.36,0.58);
\draw (1.36,0.38)--(2.18,0.38);
\draw (1.77,0.38)--(1.77,1.00)--(0.98,1.00);
\draw[-{book tip}] (0.98,1.00)--(0.98,0.60);
\draw (2.18,-0.04)--(2.18,0.80);
\draw (2.35,0.16)--(2.35,0.60);
\node at (1.88,-0.05) {$+$};
\node at (2.62,-0.05) {$-$};
\end{tikzpicture}
\end{center}
}{
\smallskip\noindent\textbf{Решение от GPT}\par
При перемещении ползунка влево сопротивление включённой части реостата уменьшается, поэтому ток в катушке №~1 и создаваемая им индукция возрастают: утверждения 1 и 2 верны. Модуль магнитного потока через катушку №~2 увеличивается, поэтому утверждение 3 неверно. Ток от положительного полюса источника идёт по катушке №~1 по часовой стрелке, создавая поле от наблюдателя. По правилу Ленца поле индукционного тока катушки №~2 направлено к наблюдателю, а сам ток идёт против часовой стрелки: утверждения 4 и 5 неверны. Ответ: 1, 2.
}

% Source F4BE25: https://ege.fipi.ru/bank/questions.php?search=1&proj=BA1F39653304A5B041B656915DC36B38&qid=F4BE25
\BankTask{F4BE25}{
На железный сердечник надеты две катушки, как показано на рисунке. По правой катушке пропускают ток, который меняется согласно приведённому графику.
\begin{enumerate}
\item В промежутках времени $0\text{–}1\,\text{с}$ и $3\text{–}5\,\text{с}$ направления тока в левой катушке одинаковы.
\item В промежутке времени $2\text{–}3\,\text{с}$ сила тока в левой катушке равна $0$.
\item Модули силы тока в левой катушке в промежутках времени $1\text{–}2\,\text{с}$ и $3\text{–}5\,\text{с}$ одинаковы.
\item В промежутке $0\text{–}2\,\text{с}$ модуль индукции магнитного поля в сердечнике равен $0$.
\item В левой катушке сила тока в промежутке времени $0\text{–}1\,\text{с}$ по модулю больше, чем в промежутке времени $3\text{–}5\,\text{с}$.
\end{enumerate}
\begin{center}
\begin{tikzpicture}[book diagram,x=1cm,y=1cm]
\draw (0,0.8) rectangle (4.4,2.25);
\foreach \x in {0.25,0.35,...,1.75}{
\draw[line width=0.3pt] (\x,0.8) .. controls (\x-0.15,1.15) and (\x-0.10,1.90) .. (\x+0.10,2.25);
\draw[line width=0.3pt] (\x+0.10,2.25) .. controls (\x+0.24,2.55) and (\x+0.25,2.33) .. (\x+0.18,2.15);
\draw[line width=0.3pt] (\x+0.18,2.15) .. controls (\x+0.01,1.55) and (\x+0.01,1.00) .. (\x+0.10,0.8);
}
\foreach \x in {2.5,2.6,...,4.0}{
\draw[line width=0.3pt] (\x,0.8) .. controls (\x-0.15,1.15) and (\x-0.10,1.90) .. (\x+0.10,2.25);
\draw[line width=0.3pt] (\x+0.10,2.25) .. controls (\x+0.24,2.55) and (\x+0.25,2.33) .. (\x+0.18,2.15);
\draw[line width=0.3pt] (\x+0.18,2.15) .. controls (\x+0.01,1.55) and (\x+0.01,1.00) .. (\x+0.10,0.8);
}
\draw (0.25,0.8)--(0.25,0.05);
\draw (0.25,-0.12) rectangle (0.98,0.22);
\draw (0.98,0.05)--(1.16,0.05);
\draw (1.49,0.05) circle[radius=0.33];
\node at (1.49,0.05) {$\mathrm{A}$};
\draw (1.82,0.05)--(1.90,0.05)--(1.90,0.8);
\draw (2.5,0.8)--(2.5,0.05)--(3.04,0.05);
\draw (3.12,0.05) circle[radius=0.065];
\draw (3.62,0.05) circle[radius=0.065];
\draw (3.69,0.05)--(4.16,0.05)--(4.16,0.8);
\draw (3.02,-0.07)--(3.22,0.17);
\draw (3.52,-0.07)--(3.72,0.17);
\end{tikzpicture}

\begin{tikzpicture}[book diagram,x=0.65cm,y=0.325cm]
\draw[xstep=0.5,ystep=1,gray,very thin] (0,-4) grid (5,4);
\draw[-{book tip}] (0,-5.2)--(0,5.9);
\draw[-{book tip}] (0,0)--(6,0);
\node[above left] at (0,5.1) {$i,\,\text{А}$};
\node[below right] at (5.7,-0.05) {$t,\,\text{с}$};
\foreach \x in {1,2,3,4,5}{\node[below,fill=white,inner sep=1pt] at (\x,0) {$\x$};}
\foreach \y in {-4,-2,0,2,4}{\node[left] at (0,\y) {$\y$};}
\draw[thick] (0,-4)--(2,4)--(3,4)--(5,0);
\end{tikzpicture}
\end{center}
}{
\smallskip\noindent\textbf{Решение от GPT}\par
Индукционный ток в левой катушке определяется изменением потока: $\mathcal E=-M\,\Delta i/\Delta t$, поэтому его модуль пропорционален модулю наклона графика. На участках $0$–$2\,\text{с}$ и $3$–$5\,\text{с}$ наклоны равны соответственно $4\,\text{А/с}$ и $-2\,\text{А/с}$: направления индукционного тока противоположны, а модули отличаются вдвое, поэтому утверждения 1 и 3 неверны, 5 верно. На участке $2$–$3\,\text{с}$ ток правой катушки постоянен, индукционный ток отсутствует: утверждение 2 верно. На участке $0$–$2\,\text{с}$ магнитная индукция не равна нулю на всём промежутке, поскольку ток обращается в нуль только при $t=1\,\text{с}$: утверждение 4 неверно. Ответ: 2, 5.
}

% Source Fe341D: https://ege.fipi.ru/bank/questions.php?search=1&proj=BA1F39653304A5B041B656915DC36B38&qid=Fe341D
\BankTask{Fe341D}{
Проволочная рамка площадью $50\,\text{см}^{2}$ помещена в однородное магнитное поле так, что плоскость рамки перпендикулярна вектору индукции $\vec B$. Проекция $B_n$ индукции магнитного поля на нормаль к плоскости рамки изменяется во времени $t$ согласно графику на рисунке.
\begin{enumerate}
\item Модуль ЭДС индукции, возникающей в рамке в промежутке времени от $5$ до $8\,\text{мс}$, равен $4\,\text{В}$.
\item Направления индукционного тока в рамке в промежутках времени от $0$ до $2\,\text{мс}$ и от $5$ до $8\,\text{мс}$ совпадают.
\item Модуль изменения магнитного потока в рамке минимален в промежутке времени от $2$ до $4\,\text{мс}$.
\item Модуль силы индукционного тока в рамке максимален в промежутке времени от $0$ до $2\,\text{мс}$.
\item ЭДС индукции в рамке отлична от $0$ только в промежутках времени от $0$ до $2\,\text{мс}$ и от $5$ до $8\,\text{мс}$.
\end{enumerate}
\begin{center}
\begin{tikzpicture}[book diagram,x=0.65cm,y=0.65cm]
\draw[step=1,gray,very thin] (0,-2) grid (8,3);
\draw[-{book tip}] (0,-2.25)--(0,4.2);
\draw[-{book tip}] (0,0)--(9.15,0);
\node[above left] at (0,3.85) {$B_n,\,\text{Тл}$};
\node[below right] at (8.85,-0.05) {$t,\,\text{мс}$};
\foreach \x in {1,2,3,4,5,6,7,8}{\node[below,fill=white,inner sep=1pt] at (\x,0) {$\x$};}
\foreach \y in {-2,-1,0,1,2,3}{\node[left] at (0,\y) {$\y$};}
\draw[thick] (0,-2)--(2,2)--(4,-1)--(5,-1)--(8,3);
\end{tikzpicture}
\end{center}
}{
\smallskip\noindent\textbf{Решение от GPT}\par
По закону Фарадея $|\mathcal E|=S|\Delta B_n/\Delta t|$, где $S=5\cdot10^{-3}\,\text{м}^2$; на участках $0$–$2$, $2$–$4$, $4$–$5$ и $5$–$8\,\text{мс}$ получаем соответственно $10$, $7{,}5$, $0$ и $20/3\,\text{В}$, поэтому утверждение 1 неверно. На первом и последнем участках $B_n$ возрастает, следовательно, направления индукционного тока совпадают: утверждение 2 верно. Минимальное изменение потока равно нулю на участке $4$–$5\,\text{мс}$, поэтому утверждение 3 неверно. При постоянном сопротивлении $|I|=|\mathcal E|/R$ максимально на участке $0$–$2\,\text{мс}$: утверждение 4 верно. На участке $2$–$4\,\text{мс}$ ЭДС также отлична от нуля, поэтому утверждение 5 неверно. Ответ: 2, 4.
}

% Source c584F9: https://ege.fipi.ru/bank/questions.php?search=1&proj=BA1F39653304A5B041B656915DC36B38&qid=c584F9
\BankTask{c584F9}{
По гладким параллельным проводящим рельсам, которые расположены на диэлектрической подножке и замкнуты на резистор $R$, перемещают лёгкий тонкий проводник $MN$. Рельсы, резистор и проводник образуют контур, который находится в однородном магнитном поле с вектором магнитной индукции $\vec{B}$, как показано на рисунке а. При движении проводника площадь, ограниченная контуром, изменяется с течением времени так, как указано на рисунке б. ЭДС самоиндукции, возникающей в контуре, пренебречь.

\begin{enumerate}
\item В интервале времени от $20$ до $40\,\text{с}$ в резисторе течёт ток.
\item Модуль ЭДС индукции, возникающей в контуре, достигает максимального значения в интервале времени от $40$ до $50\,\text{с}$.
\item В интервале времени от $0$ до $20\,\text{с}$ сила Ампера, действующая на проводник, направлена влево.
\item Величина силы, прикладываемой к проводнику для его перемещения, максимальна в интервале времени от $40$ до $50\,\text{с}$.
\item Поскольку рельсы гладкие, для перемещения проводника к нему не надо прикладывать силу ни в один из интервалов времени, представленных на графике на рис. б.
\end{enumerate}

\begin{center}
\begin{tikzpicture}[book diagram,x=0.8cm,y=0.8cm]
\draw[line width=2pt] (0,0)--(6,0) (0,1.95)--(6,1.95);
\draw (1.2,-0.4)--(1.2,2.35);
\node[above left] at (1.2,2.05) {$N$};
\node[below left] at (1.2,-0.05) {$M$};
\draw (4.8,1.95)--(4.8,1.48) (4.8,0.5)--(4.8,0);
\draw (4.53,0.5) rectangle (5.07,1.48);
\node[right] at (5.1,0.99) {$R$};
\draw (2.96,0.94) circle (0.33);
\draw (2.73,0.71)--(3.19,1.17) (2.73,1.17)--(3.19,0.71);
\node[below right] at (3.32,0.93) {$\vec B$};
\node at (3,-0.8) {Рис. а};
\end{tikzpicture}
\qquad
\begin{tikzpicture}[book diagram,x=0.052cm,y=2.7cm]
\draw[-{book tip}] (0,0)--(65,0) node[right] {$t,\,\text{с}$};
\draw[-{book tip}] (0,0)--(0,1.04) node[above left] {$S,\,\text{м}^2$};
\node[below left] at (0,0) {$0$};
\foreach \x in {10,20,30,40,50} {
\draw (\x,-0.025)--(\x,0.025);
\node[below] at (\x,-0.035) {$\x$};
}
\foreach \x in {5,15,25,35,45,55} {\draw (\x,-0.014)--(\x,0.014);}
\foreach \y/\lab in {0.2/{0{,}2},0.4/{0{,}4},0.6/{0{,}6},0.8/{0{,}8}} {
\draw (-1.3,\y)--(1.3,\y);
\node[left] at (-1.5,\y) {$\lab$};
}
\foreach \y in {0.1,0.3,0.5,0.7,0.9} {\draw (-0.7,\y)--(0.7,\y);}
\draw[thick] (0,0.2)--(20,0.8)--(40,0.8)--(50,0.2);
\draw (0,0.2)--(18,0.2) (20,0.1)--(20,0.36) (37,0.2)--(50,0.2) (40,0.44)--(40,0.8);
\node at (30,-0.31) {Рис. б};
\end{tikzpicture}
\end{center}
}{
\smallskip\noindent\textbf{Решение от GPT}\par
По закону Фарадея $|\mathcal E|=B|\Delta S/\Delta t|$: на интервалах $0$–$20$, $20$–$40$ и $40$–$50\,\text{с}$ модули скорости изменения площади равны соответственно $0{,}03$, $0$ и $0{,}06\,\text{м}^2/\text{с}$, поэтому утверждение 1 неверно, а 2 верно. В первые $20\,\text{с}$ площадь увеличивается, проводник движется влево, а сила Ампера по правилу Ленца направлена вправо: утверждение 3 неверно. При пренебрежении инерцией лёгкого проводника прикладываемая сила уравновешивает силу Ампера, $F=BIl=\frac{B^2l}{R}|dS/dt|$, поэтому утверждение 4 верно. Гладкость рельсов устраняет трение, но не магнитное торможение, поэтому утверждение 5 неверно. Ответ: 2, 4.
}

% Source 9CFC20: https://ege.fipi.ru/bank/questions.php?search=1&proj=BA1F39653304A5B041B656915DC36B38&qid=9CFC20
\BankTask{9CFC20}{
В катушке индуктивностью $20\,\text{мГн}$ сила тока $I$ зависит от времени $t$, как показано на графике, приведённом на рисунке.

\begin{enumerate}
\item Модуль ЭДС самоиндукции в катушке максимален в интервале времени от $0$ до $1\,\text{с}$.
\item Энергия магнитного поля катушки в интервале времени от $2$ до $3\,\text{с}$ равна $40\,\text{мДж}$.
\item Модуль ЭДС самоиндукции в катушке в интервале времени от $3$ до $5\,\text{с}$ равен $10\,\text{мВ}$.
\item Модуль ЭДС самоиндукции в катушке минимален в интервале времени от $3$ до $4\,\text{с}$.
\item Модуль скорости изменения тока в катушке максимален в интервале времени от $5$ до $6\,\text{с}$.
\end{enumerate}

\begin{center}
\begin{tikzpicture}[book diagram,x=0.75cm,y=0.75cm]
\draw[thin] (0,1)--(6,1) (0,2)--(6,2) (0,3)--(6,3) (0,4)--(6,4);
\foreach \x in {1,2,3,4,5,6} {\draw[thin] (\x,0)--(\x,4);}
\draw[-{book tip}] (0,0)--(6.75,0) node[right] {$t,\,\text{с}$};
\draw[-{book tip}] (0,0)--(0,4.9) node[above left] {$I,\,\text{А}$};
\node[below left] at (0,0) {$0$};
\foreach \x in {1,2,3,4,5,6} {\node[below] at (\x,0) {$\x$};}
\foreach \y in {1,2,3,4} {\node[left] at (0,\y) {$\y$};}
\draw[thick] (0,0)--(2,4)--(3,4)--(5,3)--(6,0);
\end{tikzpicture}
\end{center}
}{
\smallskip\noindent\textbf{Решение от GPT}\par
Используем $|\mathcal E|=L|\Delta I/\Delta t|$: на интервалах $0$–$2$, $2$–$3$, $3$–$5$ и $5$–$6\,\text{с}$ модуль скорости изменения тока равен соответственно $2$, $0$, $0{,}5$ и $3\,\text{А}/\text{с}$, поэтому утверждение 1 неверно, а 5 верно. При $2<t<3\,\text{с}$ энергия $W=LI^2/2=0{,}02\cdot4^2/2=0{,}16\,\text{Дж}=160\,\text{мДж}$, поэтому утверждение 2 неверно. На интервале $3$–$5\,\text{с}$ получаем $|\mathcal E|=0{,}02\cdot0{,}5=0{,}01\,\text{В}=10\,\text{мВ}$, поэтому утверждение 3 верно. Минимальная ЭДС равна нулю на интервале $2$–$3\,\text{с}$, поэтому утверждение 4 неверно. Ответ: 3, 5.
}

% Source Fc855A: https://ege.fipi.ru/bank/questions.php?search=1&proj=BA1F39653304A5B041B656915DC36B38&qid=Fc855A
\BankTask{Fc855A}{
В идеальном колебательном контуре, состоящем из конденсатора и катушки индуктивности, происходят свободные электромагнитные колебания. Изменение заряда конденсатора в колебательном контуре с течением времени показано в таблице.

\begin{center}
\begin{tabular}{|c|c|c|c|c|c|c|c|c|c|c|}
\hline
$t,\ 10^{-6}\,\text{с}$ & $0$ & $1$ & $2$ & $3$ & $4$ & $5$ & $6$ & $7$ & $8$ & $9$ \\
\hline
$q,\ 10^{-9}\,\text{Кл}$ & $1$ & $0{,}71$ & $0$ & $-0{,}71$ & $-1$ & $-0{,}71$ & $0$ & $0{,}71$ & $1$ & $0{,}71$ \\
\hline
\end{tabular}
\end{center}

\begin{enumerate}
\item Период колебаний равен $8\times10^{-6}\,\text{с}$.
\item Частота колебаний равна $250\,\text{кГц}$.
\item В момент времени $t=2\times10^{-6}\,\text{с}$ модуль силы тока в катушке индуктивности максимален.
\item В момент времени $t=8\times10^{-6}\,\text{с}$ энергия магнитного поля катушки индуктивности максимальна.
\item В момент времени $t=4\times10^{-6}\,\text{с}$ энергия электрического поля конденсатора минимальна.
\end{enumerate}
}{
\smallskip\noindent\textbf{Решение от GPT}\par
Заряд возвращается к исходному максимуму через $T=8\times10^{-6}\,\text{с}$, поэтому утверждение 1 верно. Частота $\nu=1/T=125\,\text{кГц}$, поэтому утверждение 2 неверно. При $t=2\times10^{-6}\,\text{с}$ заряд равен нулю, а модуль тока максимален, поэтому утверждение 3 верно. При $t=8\times10^{-6}\,\text{с}$ заряд максимален, ток и энергия магнитного поля равны нулю, поэтому утверждение 4 неверно. При $t=4\times10^{-6}\,\text{с}$ модуль заряда максимален и энергия $W_{\text{эл}}=q^2/(2C)$ максимальна, поэтому утверждение 5 неверно. Ответ: 1, 3.
}

% Source E51107: https://ege.fipi.ru/bank/questions.php?search=1&proj=BA1F39653304A5B041B656915DC36B38&qid=E51107
\BankTask{E51107}{
На рисунке приведён график зависимости силы тока от времени в колебательном контуре, образованном конденсатором и катушкой, индуктивность которой равна $0{,}3\,\text{Гн}$.
\begin{enumerate}
\item Период электромагнитных колебаний в контуре равен $2\,\text{мс}$.
\item Максимальное значение энергии электрического поля конденсатора равно $5{,}4\,\text{мкДж}$.
\item В момент времени $5\,\text{мс}$ заряд конденсатора равен нулю.
\item В момент времени $3\,\text{мс}$ энергия магнитного поля катушки достигает своего минимума.
\item За первые $6\,\text{мс}$ энергия магнитного поля катушки достигла своего максимума $2$ раза.
\end{enumerate}
\begin{center}
\begin{tikzpicture}[book diagram,x=0.62cm,y=0.14cm]
\draw[-{book tip}] (0,-8.4)--(0,10.8);
\draw[-{book tip}] (0,0)--(7.4,0);
\node[above left] at (0,9.3) {$i,\,\text{мА}$};
\node[below right] at (6.65,-0.5) {$t,\,\text{мс}$};
\foreach \x in {1,2,3,4,5,6}{\draw (\x,-0.5)--(\x,0.5);\node[below] at (\x,-0.7) {$\x$};}
\foreach \y in {-6,6}{\draw (-0.12,\y)--(0.12,\y);\node[left] at (-0.18,\y) {$\y$};}
\node[below left] at (-0.08,0) {$0$};
\draw[very thick,smooth,samples=160,domain=0:6] plot (\x,{6*sin(90*\x)});
\end{tikzpicture}
\end{center}
}{
\smallskip\noindent\textbf{Решение от GPT}\par
По графику $T=4\,\text{мс}$, поэтому утверждение 1 неверно. Максимальная энергия электрического поля равна полной энергии контура: $W_{\max}=LI_{\max}^2/2=0{,}3\cdot(6\cdot10^{-3})^2/2=5{,}4\cdot10^{-6}\,\text{Дж}$, поэтому утверждение 2 верно. При $t=5\,\text{мс}$ ток максимален, вся энергия заключена в магнитном поле, а заряд конденсатора равен нулю: утверждение 3 верно. При $t=3\,\text{мс}$ модуль тока максимален, поэтому магнитная энергия максимальна и утверждение 4 неверно. За первые $6\,\text{мс}$ её максимумы достигаются трижды: при $t=1$, $3$ и $5\,\text{мс}$, поэтому утверждение 5 неверно. Ответ: 2, 3.
}

% Source 209C4E: https://ege.fipi.ru/bank/questions.php?search=1&proj=BA1F39653304A5B041B656915DC36B38&qid=209C4E
\BankTask{209C4E}{
Школьник, изучая законы геометрической оптики, провёл опыт по преломлению света (см. рисунок). Для этого он направил узкий пучок света на стеклянную пластину.
\begin{center}
\begin{tabular}{|c|c|c|c|c|}
\hline
Угол $\alpha$ & $20^\circ$ & $40^\circ$ & $50^\circ$ & $70^\circ$ \\
\hline
$\sin\alpha$ & $0{,}34$ & $0{,}64$ & $0{,}78$ & $0{,}94$ \\
\hline
\end{tabular}
\end{center}
\begin{enumerate}
\item Угол падения равен $70^\circ$.
\item Показатель преломления стекла примерно равен $1{,}47$.
\item Угол преломления равен $50^\circ$.
\item Наблюдается полное внутреннее отражение.
\item Угол отражения равен $20^\circ$.
\end{enumerate}
\begin{center}
\begin{tikzpicture}[book diagram,scale=1.15]
\fill[gray!35] (0,0) circle (2.1);
\shade[top color=gray!60,bottom color=white] (-1.56,0) arc[start angle=180,end angle=360,radius=1.56]--cycle;
\draw (-1.56,0)--(1.56,0) arc[start angle=0,end angle=-180,radius=1.56];
\draw (0,0) circle (2.1);
\foreach \a in {0,5,...,355}{\draw (\a:1.98)--(\a:2.1);}
\foreach \a in {0,10,...,350}{\draw (\a:1.88)--(\a:2.1);}
\foreach \a in {0,30,...,330}{\draw (\a:1.81)--(\a:2.1);}
\node[rotate=90] at (0:1.72) {$0$};
\node[rotate=60] at (30:1.72) {$30$};
\node[rotate=30] at (60:1.72) {$60$};
\node at (90:1.72) {$90$};
\node[rotate=-30] at (120:1.72) {$60$};
\node[rotate=-60] at (150:1.72) {$30$};
\node[rotate=90] at (180:1.72) {$0$};
\node[rotate=60] at (210:1.72) {$30$};
\node[rotate=30] at (240:1.72) {$60$};
\node[rotate=180] at (270:1.72) {$90$};
\node[rotate=-30] at (300:1.72) {$60$};
\node[rotate=-60] at (330:1.72) {$30$};
\draw[thick] (0,0)--(0,1.56);
\draw[gray!80,very thick,-{book tip}] (160:2.9)--(160:2.13);
\draw[white,line width=0.9pt] (160:2.1)--(0,0)--(20:2.1);
\draw[gray!80,very thick,-{book tip}] (20:2.1)--(20:2.9);
\draw[white,line width=1.5pt] (0,0)--(-50:2.1);
\draw[gray!80,line width=0.5pt] (0,0)--(-50:2.1);
\draw[gray!80,very thick,-{book tip}] (-50:2.1)--(-50:2.9);
\end{tikzpicture}
\end{center}
}{
\smallskip\noindent\textbf{Решение от GPT}\par
Угол падения отсчитывают от нормали: $\alpha=90^\circ-20^\circ=70^\circ$, поэтому утверждение 1 верно. Угол преломления $\beta=90^\circ-50^\circ=40^\circ$, а показатель преломления стекла $n=\sin70^\circ/\sin40^\circ=0{,}94/0{,}64\approx1{,}47$: утверждение 2 верно, утверждение 3 неверно. Преломлённый луч присутствует, поэтому полного внутреннего отражения нет и утверждение 4 неверно. Угол отражения равен углу падения, то есть $70^\circ$, поэтому утверждение 5 неверно. Ответ: 1, 2.
}

\end{document}
